\documentclass[]{vvsu}

\vvsuyear{2025}

%%%%%%%%%%%%%%%%%%%

\usepackage{graphicx} % для изображений
\usepackage{tabularray} % для таблиц
\usepackage{siunitx} % для обозначений (процент, градус)
\usepackage{listings} % для листингов кода

% Список путей, где будут искаться изображения и файлы
\graphicspath{{images/}}

% Автор документа
\author{Т. Г. Пономаренко}

% Настройка стилей для листингов кода
\input{listing_styles.tex}

%%%%%%%%%%%%%%%%%%%

\begin{document}

% Шапка
\vvsuhead{\linespread{1}\selectfont{}МИНОБРНАУКИ РОССИИ\\
\vspace{10pt}Федеральное государственное бюджетное образовательное учреждение\\
высшего образования\\
\fontsize{13}{13}\selectfont{}<<ВЛАДИВОСТОКСКИЙ ГОСУДАРСТВЕННЫЙ УНИВЕРСИТЕТ>>\\
(ФГБОУ ВО <<ВВГУ>>)\\
\vspace{10pt}\fontsize{12}{12}\selectfont{}ИНСТИТУТ ИНФОРМАЦИОННЫХ ТЕХНОЛОГИЙ И АНАЛИЗА ДАННЫХ\\
КАФЕДРА ИНФОРМАЦИОННЫХ ТЕХНОЛОГИЙ И СИСТЕМ}

% Название отчета
\title{Отчет\\по лабораторной работе №6}
\subtitle{по дисциплине\\<<Информатика и программирование>>}

% Участники работы
\member{Студент\\ гр. БИН-25-3}{Т. Г. Пономаренко}
\member{Ассистент\\ преподавателя}{М.В. Водяницкий}

% Вывод титульника
\maketitle

% Задание
\begin{addition}{Задание}
  Выполнить задания и оформить отчет по стандартам ВВГУ.

 \textit{\textbf{Задание 1.}}  
  Написать функцию, которая конвертирует время из одной величины в другую.  
  На вход подаются:  
  \begin{vvsu_itemize}
    \item число (величина времени)
    \item исходная единица измерения
    \item единица измерения, в которую нужно перевести
  \end{vvsu_itemize}
  Функция должна вернуть конвертированное значение.

  Пример:\\
    Введите время для конвертации: 4 h m\\
    Результат: 240m

  \textit{\textbf{Задание 2.}}  
  Пользователь делает вклад в банке в размере \texttt{a} рублей сроком на \texttt{n} лет. Процент по вкладу зависит от суммы и срока.

  \textbf{Зависимость от суммы:}
  \begin{vvsu_itemize}
    \item каждые 10\,000 рублей увеличивают ставку на 0.3\%
    \item суммарное увеличение не может превышать 5\%
    \item минимальный вклад — 30\,000 рублей
  \end{vvsu_itemize}

  \textbf{Зависимость от срока:}
  \begin{vvsu_itemize}
    \item первые 3 года — 3\%
    \item от 4 до 6 лет — 5\%
    \item более 6 лет — 2\%
  \end{vvsu_itemize}

  Необходимо написать функцию, которая рассчитывает \textbf{прибыль пользователя без учета первоначально вложенной суммы}. Используется сложный процент.

  Пример:\\
    Ввод: 30000 3\\
    Прибыль: 3648.67 руб.

  \textit{\textbf{Задание 3.}}  
  Написать функцию для вывода всех простых чисел в заданном диапазоне. Нужно учитывать некорректные данные (например, начало больше конца или диапазон без простых чисел).

  На вход подаются два числа: начало и конец диапазона (включительно).  
  На выходе — список всех простых чисел или сообщение об ошибке.

  Пример:\\
    Начало диапазона: 1\\
    Конец диапазона: 10\\
    2 3 5 7

  Пример:\\
    Начало диапазона: 0\\
    Конец диапазона: 1\\
    Error!

  \textit{\textbf{Задание 4.}}  
  Реализовать функцию сложения двух матриц.  
  При сложении двух матриц получается новая матрица того же размера, где каждый элемент — это сумма элементов с тем же индексом из двух исходных матриц.

  \textbf{Ограничения:}
  \begin{vvsu_itemize}
    \item складывать можно только матрицы одинакового размера
    \item размер матрицы должен быть строго больше 2 (например, 3×3, 4×4 и т.д.)
    \item при нарушении условий нужно вывести сообщение об ошибке
  \end{vvsu_itemize}

  На вход подаются:
  \begin{vvsu_list}
    \item размер матрицы \texttt{n} (для квадратной матрицы \texttt{n × n})
    \item элементы первой матрицы (по строкам, через пробел)
    \item элементы второй матрицы (в таком же формате)
  \end{vvsu_list}

  Результат — новая матрица (в том же формате), либо сообщение об ошибке.

  \textit{\textbf{Задание 5.}}  
  Написать функцию, которая определяет, является ли строка палиндромом.  
  Палиндром — это строка, которая читается одинаково слева направо и справа налево (без учета пробелов, регистра и знаков препинания).

  На вход подается строка.  
  На выходе:
  \begin{vvsu_itemize}
    \item \texttt{Да}, если это палиндром
    \item \texttt{Нет}, если это не палиндром
  \end{vvsu_itemize}

  Пример:\\
    А роза упала на лапу Азора\\
    Да

  Пример:\\
    Алфавитный порядок\\
    Нет
\end{addition}


% Содержание
\toc

% Глава - Выполнение работы
\section{Выполнение работы}

% Подглава - Задание 1
\subsection{Задание 1}

1. Объявление функции - Создаём функцию convert\_time с тремя параметрами
2. Проверка входной единицы "часы" - Если исходная единица равна "h"
3. Расчёт секунд из часов - Умножаем число на 3600 (секунд в одном часе)  
4. Проверка входной единицы "минуты" - Если исходная единица равна `"m"`  
5. Расчёт секунд из минут - Умножаем число на 60 (секунд в одной минуте)  
6. Случай для секунд - Если исходная единица равна `"s"`, оставляем без изменений  
7. Другой случай - Если единица не распознана, считаем что это секунды  
8. Проверка выходной единицы "часы" - Если целевая единица равна `"h"`  
9. Расчёт часов из секунд - Делим общее количество секунд на 3600  
10. Проверка выходной единицы "минуты" - Если целевая единица равна `"m"`  
11. Расчёт минут из секунд - Делим общее количество секунд на 60  
12. Случай для секунд на выходе - Если целевая единица равна `"s"`, оставляем без изменений  
13. Другой случай - Если единица не распознана, оставляем секунды  
14. Вывод результата - Преобразуем число в строку и добавляем целевую единицу измерения.На рисунке \ref{fig:code_task_1} представлен код программы.

\begin{vvsu_figure}{Листинг программы для задания 1}{fig:code_task_1}
  \begin{minipage}{.75\textwidth}
    \lstinputlisting[language=Python,basicstyle=\fontsize{10}{10}\linespread{1}\selectfont\ttfamily]{code/task1.py}
  \end{minipage}
\end{vvsu_figure}

% Подглава - Задание 2
\subsection{Задание 2}

1. Объявление функции - Создаём функцию `vklad` с двумя параметрами (сумма и срок)  
2. Проверка минимальной суммы - Проверяем, что сумма больше или равна 30000  
3. Вывод ошибки при малой сумме - Если сумма меньше 30000, выводим "Error!" и завершаем  
4. Расчёт бонусной ставки - За каждые 10000 добавляем 0.3% к ставке  
5. Ограничение бонусной ставки - Максимальный бонус не может превышать 5%  
6. Определение базовой ставки по сроку - Для срока до 3 лет: 3%, от 4 до 6 лет: 5%, больше 6 лет: 2%  
7. Расчёт общей процентной ставки - Складываем базовую и бонусную ставки  
8. Инициализация переменной для накопления - Начинаем с исходной суммы  
9. Инициализация счётчика лет - Устанавливаем начальное значение 0  
10. Цикл начисления процентов - Повторяем для каждого года вклада  
11. Начисление процентов за год - Увеличиваем сумму на процент от текущей суммы  
12. Увеличение счётчика лет - Переходим к следующему году  
13. Расчёт прибыли - Вычитаем из итоговой суммы первоначальный вклад  
14. Вывод результата с округлением - Округляем прибыль до двух знаков после запятой На рисунке \ref{fig:code_task_2} представлен код программы.

\begin{vvsu_figure}{Листинг программы для задания 2}{fig:code_task_2}
  \begin{minipage}{.75\textwidth}
    \lstinputlisting[language=Python,basicstyle=\fontsize{10}{10}\linespread{1}\selectfont\ttfamily]{code/task2.py}
  \end{minipage}
\end{vvsu_figure}

% Подглава - Задание 3
\subsection{Задание 3}

1. Объявление функции - Создаём функцию `primes` с двумя параметрами (начало и конец диапазона)  
2. Создание пустого списка - Инициализируем список для хранения простых чисел  
3. Цикл по числам диапазона - Перебираем числа от начала до конца включительно, начиная минимум с 2  
4. Проверка числа на простоту - Рассматриваем только числа больше 1  
5. Поиск делителей - Проверяем делители числа от 2 до квадратного корня из числа  
6. Проверка делимости - Если число делится на какой-либо делитель без остатка  
7. Выход из внутреннего цикла - Прерываем проверку, так как число составное  
8. Добавление простого числа - Если цикл завершился без прерывания, добавляем число в список  
9. Возврат результата - Возвращаем список простых чисел или "Error!" если список пуст  
10. Ввод начала диапазона - Запрашиваем у пользователя начало диапазона  
11. Ввод конца диапазона - Запрашиваем у пользователя конец диапазона  
12. Вывод результата - Выводим все простые числа через пробел На рисунке \ref{fig:code_task_3} представлен код программы.

\begin{vvsu_figure}{Листинг программы для задания 3}{fig:code_task_3}
  \begin{minipage}{.75\textwidth}
    \lstinputlisting[language=Python,basicstyle=\fontsize{10}{10}\linespread{1}\selectfont\ttfamily]{code/task3.py}
  \end{minipage}
\end{vvsu_figure}

% Подглава - Задание 4
\subsection{Задание 4}

1. Объявление функции - Создаём функцию   без параметров  
2. Начало блока обработки ошибок - Помещаем код в конструкцию `try` для отлова исключений  
3. Ввод размера матрицы - Запрашиваем у пользователя размер квадратной матрицы  
4. Проверка минимального размера - Проверяем, что размер больше 2  
5. Вывод ошибки при малом размере - Если размер меньше или равен 2, выводим "Error!" и завершаем  
6. Ввод первой матрицы - Запрашиваем элементы первой матрицы построчно  
7. Преобразование строки в список чисел - Разбиваем строку на числа и создаём список  
8. Создание матрицы 1 - Формируем список списков (матрицу) из введённых строк  
9. Ввод второй матрицы - Запрашиваем элементы второй матрицы построчно  
10. Создание матрицы 2 - Формируем вторую матрицу аналогичным образом  
11. Вывод заголовка результата - Печатаем "Результат:" перед выводом суммы матриц  
12. Цикл по строкам матрицы - Перебираем строки от 0 до n-1  
13. Сложение элементов матриц - Для каждого столбца складываем элементы из обеих матриц  
14. Формирование строки результата - Преобразуем суммы в строки и соединяем пробелами  
15. Вывод строки результата - Печатаем каждую строку результирующей матрицы  
16. Обработка ошибок ввода - Если возникает любая ошибка, переходим в блок `except`  
17. Вывод сообщения об ошибке - При любой ошибке выводим "Error!"  
18. Вызов функции - Запускаем выполнение функции На рисунке \ref{fig:code_task_4} представлен код решения.

\begin{vvsu_figure}{Листинг программы для задания 4}{fig:code_task_4}
  \begin{minipage}{.75\textwidth}
    \lstinputlisting[language=Python,basicstyle=\fontsize{10}{10}\linespread{1}\selectfont\ttfamily]{code/task4.py}
  \end{minipage}
\end{vvsu_figure}


% Подглава - Задание 5
\subsection{Задание 5}

1. Объявление функции - Создаём функцию  без параметров  
2. Ввод строки - Запрашиваем строку и удаляем пробелы по краям  
3. Создание пустой строки - Инициализируем строку для очищенного текста  
4. Очистка строки - Убираем все символы кроме букв и цифр, переводим в нижний регистр  
5. Создание перевёрнутой строки - Получаем строку в обратном порядке  
6. Сравнение строк - Проверяем, равна ли строка своему перевёрнутому варианту  
7. Вывод результата - Если строки равны: "Да", иначе: "Нет"  
8. Запуск функции - Вызываем функцию для выполнения На рисунке \ref{fig:code_task_5} представлен код программы.

\begin{vvsu_figure}{Листинг программы для задания 5}{fig:code_task_5}
  \begin{minipage}{.75\textwidth}
    \lstinputlisting[language=Python,basicstyle=\fontsize{10}{10}\linespread{1}\selectfont\ttfamily]{code/task5.py}
  \end{minipage}
\end{vvsu_figure}

\end{document}